% ==============================================================================
% frontmatter.tex -- Front matter for Arithmetic Moduli of Elliptic Curves
% ==============================================================================

% Half-title page
\begin{titlepage}
  \vspace*{4cm}
  \begin{center}
    {\large\scshape Annals of Mathematics Studies\par}
    \vspace{0.3cm}
    {\normalsize Number 108\par}
    \vspace{2.5cm}
    {\Large\bfseries ARITHMETIC MODULI OF ELLIPTIC CURVES\par}
  \end{center}
  \clearpage

  % Title page
  \vspace*{2cm}
  \begin{center}
    {\Large\bfseries ARITHMETIC MODULI\par}
    \vspace{0.3cm}
    {\Large\bfseries OF ELLIPTIC CURVES\par}
    \vspace{2.5cm}
    {\large\scshape By\par}
    \vspace{0.5cm}
    {\Large\bfseries Nicholas M. Katz and Barry Mazur\par}
    \vfill
    {\large\scshape Princeton University Press\par}
    \vspace{0.2cm}
    {\normalsize Princeton, New Jersey\par}
    \vspace{0.2cm}
    {\normalsize 1985\par}
  \end{center}
  \clearpage

  % Copyright page
  \vspace*{\fill}
  \begin{flushleft}
    \small
    Copyright \copyright\ 1985 by Princeton University Press\\
    \textsc{All Rights Reserved}\\[1em]
    The Annals of Mathematics Studies are edited by\\
    William Browder, Robert P. Langlands, John Milnor, and Elias M. Stein\\[1em]
    Corresponding editors:\\
    Stefan Hildebrandt, H. Blaine Lawson, Louis Nirenberg, and David Vogan\\[1.5em]
    ISBN 0-691-08349-5 (cloth)\\
    ISBN 0-691-08352-5 (paper)\\[1em]
    Printed in the United States of America\\
    by Princeton University Press, 41 William Street\\
    Princeton, New Jersey
  \end{flushleft}
\end{titlepage}

% Table of Contents
\tableofcontents
\cleardoublepage

% Introduction
\chapter*{Introduction}
\addcontentsline{toc}{chapter}{Introduction}
\markboth{Introduction}{Introduction}

This book is devoted to giving an account of the arithmetic theory of the moduli spaces of elliptic curves. The main emphasis is on understanding the behavior of these moduli spaces at primes dividing the ``level'' of the moduli problem being considered. Until recently, this seemed a very difficult problem, because one had no a priori construction of these spaces at the ``bad'' primes. One defined them as schemes over, say, $\mathbb{Z}[1/N]$ as the solution to some well-posed moduli problem which only made sense for elliptic curves over rings in which $N$ was invertible, and then one used a process of normalization to extend them to schemes over $\mathbb{Z}$, e.g., one took the ``$\operatorname{Proj}$'' of the graded subring of the ring of all modular forms of the type in question consisting of those with integral Fourier (``$q$-expansion'') coefficients at the cusps. This procedure produced a scheme over $\mathbb{Z}$, but one had no idea of what moduli interpretation this scheme had, nor a fortiori did one have any idea of the modular interpretation of its reduction modulo $p$, for $p$ a prime dividing the level.

Historically, the only case where this question was in any way satisfactorily understood was the case of $\Gamma_0(p)$, which made universal sense as the moduli problem ``$p$-isogenies'', or ``finite flat subgroup-schemes of rank $p$.'' This modular interpretation was implicitly known to Kronecker, for whom it appears as the statement that the ``modular equation of degree $p$, reduced mod $p$, is the curve in the $(j_1, j_2)$-plane
\[
(j_1 - j_2^p)(j_1^p - j_2) = 0.''
\]
One knows the crucial role that the reinterpretation of this Kronecker congruence by Eichler-Shimura as the ``congruence relation''
\[
\mathbf{T}_p = \mathbf{F} + \mathbf{V} \pmod p
\]
played in the reduction of the Ramanujan conjecture to Weil's ``Riemann Hypothesis for varieties over finite fields.''

Eichler-Shimura made use of this relation to prove that for all but finitely many $p$ the Ramanujan conjecture held for any given cusp form of weight two and level $N$ which is a simultaneous eigenfunction of all $T_p$ with $p$ not dividing $N$. The method was intrinsically incapable of specifying the exceptional $p$, which were believed to consist only of primes dividing $N$.

Partly in order to settle definitively this question of exceptional $p$ for weight two forms, Igusa, in a brilliant series of papers (\cite{Ig2}, \cite{Ig3}, \cite{Ig4}, \cite{Ig5}), gave a complete and definitive account of the level $N$ moduli scheme over $\mathbb{Z}[1/N]$. Except for a difference of mathematical language, and the modular interpretation of the cusps by Deligne-Rapoport, there have been no ``improvements'' to Igusa's account of what happens over $\mathbb{Z}[1/N]$. Although Igusa's papers contain many stimulating speculations about the situation mod $p$ for ``bad'' $p$ (e.g., the footnote (\cite{Ig2}, p.~472) where he points out that the genus of $\Gamma_0(p)$ is closely related to the number of [super] singular points in characteristic $p$), there was to be no significant progress in understanding the situation at ``bad'' $p$ for another decade.

In 1968, Deligne completed the general reduction of Ramanujan's conjecture, for forms of arbitrary weight, and in particular for $\Delta$, to Weil's Riemann Hypothesis for varieties over finite fields. In his article \cite{De1}, he mentions that in fact the $\Gamma_0(p)$ moduli scheme (with suitable auxiliary prime-to-$p$ rigidification) is actually a regular scheme, and in a letter of July 10, 1970 to Parshin he proves this regularity by checking what happens at the supersingular points in characteristic $p$.

Simultaneously, another theme was developing. Shimura conjectured and Casselman \cite{Cmn1} proved that for $p = 29, 53, 61, 73, 89, 97$, the Jacobian of (the modular curve for) $\Gamma_1(p)$, modulo the Jacobian of $\Gamma_0(p)$, acquired good reduction over the field $\mathbb{Q}(\zeta_p)$. Casselman \cite{Cmn2} explained that such theorems of good reduction could be predicted by the ``Langlands philosophy'', which related, conjecturally, such questions to questions in representation theory which were already well-understood.

The paper of Deligne-Rapoport \cite{DeRa} in 1972 provided an exhaustive account of what was then known about arithmetic moduli of elliptic curves. It gave a complete account of level $N$ moduli problems over $\mathbb{Z}[1/N]$, including a modular interpretation of the compactified moduli scheme (i.e., including the cusps) as the moduli space of ``generalized elliptic curves with auxiliary structure.'' It also gave a modular interpretation over $\mathbb{Z}$ to the $\Gamma_1(p)$ moduli problem, and with it a proof that the good reduction phenomenon of Shimura-Casselman held for any $p$. Another innovation was the systematic use of algebraic stacks, as developed by Mumford \cite{Mum1} and Deligne-Mumford \cite{DeMu}.

The next significant progress came in 1974, with Drinfeld's introduction (in the context of his theory of ``elliptic modules'') of the notion of a ``full level $N$-structure'' on an elliptic curve over an arbitrary scheme, where $N$ need not be invertible, as a pair of points $P, Q$ of order $N$ such that the group-scheme $E[N]$ of points of order $N$ is equal to the sum
\[
\sum_{a,b \bmod N} [aP + bQ]
\]
as a Cartier divisor inside $E$. Drinfeld showed that with this definition, the corresponding full level $N$ moduli problem for his ``elliptic modules'' was regular. It was clear to the experts, although never published, that with Drinfeld's definition applied to usual elliptic curves, one obtained a moduli problem over $\mathbb{Z}$ which was regular, and which, over $\mathbb{Z}[1/N]$, coincided with the usual ``full level $N$'' moduli problem. In particular, one now had a modular interpretation of its reduction modulo any $p$, as the moduli space of elliptic curves, together with Drinfeld level $N$ structures, over $\mathbb{F}_p$-algebras. With this modular interpretation, it became a pleasant exercise to calculate explicitly the reduction modulo any prime $p$.

In a letter to Drinfeld of January 21, 1975, Deligne explained how the Drinfeld idea of using Cartier divisors allowed one to define universally the $\Gamma_1(N)$ problem as well as a ``balanced'' version of it by saying that a point $P$ in an elliptic curve has ``exact order $N$'' if it is killed by $N$ and if the Cartier divisor inside $E$ defined by
\[
\sum_{a \bmod N} [aP]
\]
is actually a subgroup-scheme of $E$. Deligne also explained that the resulting moduli problem was regular.

In June 1979, the present authors rediscovered Deligne's $\Gamma_1(N)$ idea, and they formulated a Drinfeldian version of $\Gamma_0(N)$ by defining a finite locally free subgroup-scheme $G$ of rank $N$ inside an elliptic curve to be cyclic if locally f.p.p.f.\ on the base one could find a point $P$ in it which generated it, in the sense that
\[
G = \sum_{a \bmod N} [aP]
\]
as Cartier divisors inside $E$. Using this definition of $\Gamma_0(N)$ as the moduli problem of ``elliptic curves together with cyclic subgroup-schemes which are finite locally free of rank $N$,'' they proved that the $\Gamma_0(N)$ problem was regular and worked out explicitly the reductions mod $p$ of all the ``standard'' moduli problems.

These calculations of special fibers, together with some intricate (due to wild ramification) calculations of the topological invariants of the special fibers, led to a direct geometric verification of a rather general theorem of good reduction which includes the Shimura-Casselman-Deligne-Rapoport theorem as a special case. For the most part, this good reduction theorem is also a consequence of the above-mentioned Langlands philosophy, which reduces it to a known question in representation theory (cf.\ \cite{La} and \cite{MW}, proof of Prop.~2, \S2, Chapter~3).

In writing this book, we have tried simultaneously to be self-contained and to be as general as possible. In the first chapter, we develop the Drinfeldian notions of level structure through the notion of equality of Cartier divisors in the ambient elliptic curve. With an eye to future applications to the moduli of abelian varieties, we formulate and prove the basic theorems in the context of an arbitrary smooth commutative group-scheme of relative dimension one. In Chapter 2, we review the theory of elliptic curves, and in Chapter 3 we introduce the four basic moduli problems $\Gamma(N)$, $\Gamma_1(N)$, $\operatorname{bal.}\Gamma_1(N)$, and $\Gamma_0(N)$. In Chapter 4, we develop the general formalism of moduli problems on the category of elliptic curves, and in Chapter 5 we prove the regularity of the four basic problems. In Chapter 6 we treat cyclicity in detail. In Chapter 7 we study the quotients of our moduli problems by finite groups, and prove their regularity.

In Chapter 8 we introduce coarse moduli schemes and compactifications, and in Chapter 9 we study moduli problems over cyclotomic integer rings. Chapter 10 is devoted to the detailed analysis of the calculus of cusps and components via the Tate curve. After a brief digression concerning ``exotic'' modular morphisms and isomorphisms in Chapter 11, the remaining three chapters are devoted to the detailed study of the degeneration at bad primes of our moduli problems: in Chapter 12 we treat Igusa curves in characteristic $p$, in Chapter 13 the reductions mod $p$ of the basic moduli problems, and in Chapter 14 the application to theorems of good reduction.

\vspace{1.5em}
\begin{flushright}
  \textsc{Nicholas M. Katz}\\
  \textsc{Barry Mazur}
\end{flushright}

\cleardoublepage
